\chapter{Conditional Rewriting}\label{chap:CTRS}
\chapterauthor{Ra\'ul Guti\'errez and Salvador Lucas}

\section{Introduction}

In \emph{conditional rewriting}~\cite{kaplan1984ConditionalRewriteRules}, the application of a (conditional) rewrite rule $\ell\to r\Leftarrow c$ to a term $s$ at a position $p$ depends not only on the (successful) matching $s|_p=\sigma(\ell)$ of the subterm $s|_p$ of $s$ at position $p$ against the left-hand side $\ell$ of the rule, but also on the conditional part $c$ which must be \emph{satisfied} by the matching substitution $\sigma$.
%
Furthermore, there are several approaches to solve such a satisfiability problem, see \cite[Chapter~7]{ohlebusch2002AdvancedTopicsTerm}.
%
\begin{example}
Consider the following  (ground) rules
\begin{eqnarray}
a & \to & b\label{ExCTRS_rule1}
\\
c & \to & d \Leftarrow b \approx a\label{ExCTRS_rule2}
\end{eqnarray}
%
Term $c$ can be rewritten with rule (\ref{ExCTRS_rule2}) depending on the \emph{interpretation} we give to $\approx$.
\begin{enumerate}
\item If $\approx$ is interpreted as \emph{rewriting in $0$ or more steps}, i.e., $\approx$ is actually $\to^*$, then $c$ \emph{cannot be rewritten} as $b$ is irreducible and the condition $b\to^*a$ cannot be satisfied.
 \item If $\approx$ is interpreted as \emph{conversion}, i.e., $\approx$ is actually $\leftrightarrow^*$, then $c$ \emph{rewrites} to $d$, as $b\leftarrow a$ holds due to rule (\ref{ExCTRS_rule1}).
\end{enumerate}
\end{example}
%
A Conditional Term Rewriting System (CTRS) ${\cal R}$ is a pair $({\cal F},R)$, where ${\cal F}$ is a signature of function symbols and $R$ is a set of (possibly) conditional rules.


\section{Syntax}

Conditional Term Rewriting Systems have the following syntax.

\begin{grammar}{CTRS Syntax}
  \begin{tabular}{lll}
  trs & ::= & \texttt{(format CTRS} cond-type\texttt{)} fun* rule* \\
  cond-type & ::= & \texttt{oriented} $|$ \texttt{join} $|$ \texttt{semi-equational} \\
  fun & ::= & \texttt{(fun} IDENTIFIER NAT\texttt{)} \\
  term & ::= & IDENTIFIER $|$ \texttt{(}IDENTIFIER term+\texttt{)} \\
  cost & ::= & \texttt{:cost} NAT \\
  rule & ::= & \texttt{(rule} term term cond* cost?\texttt{)}\\
  cond & ::= & \texttt{(=} term term\texttt{)}
  \end{tabular}
\end{grammar}

NAT and IDENTIFIER retain the same meaning as in \Cref{chap:TRS}.

Rules can appear in one of the following ways:

\begin{itemize}
  \item \texttt{(rule l r)}, where \texttt{l} and \texttt{r} are terms, for unconditional rules.
  \item \texttt{(rule l r cond1 cond2 \ldots)}, where \texttt{l} and \texttt{r} are terms and \texttt{cond}$i$ are reachability/joinability/semi-equational goals, for conditional rules.
\end{itemize}

Identifiers that are not declared as function symbols are variables.

\subsection{Restrictions}

Terms and rules must additionally satisfy the following well-formedness conditions.
\begin{enumerate}
    \item Every function symbol is declared exactly once.
    \item Function symbols are always used with their declared arity.
    \item Left-hand side of rules cannot be variable terms.
\end{enumerate}

Variables appearing in right-hand sides need not appear in the corresponding left-hand sides. At present, the termination competition supports only \texttt{oriented} CTRSs.

\section{Semantics}

In the operational description of conditional rewriting, we follow a \emph{logic-based approach} (see, e.g., \cite{plaisted1988LogicConditionalTerm,meseguer1992ConditionalRewritingLogic}), considering one-step and many-step rewritings as atoms $s\to t$ and $s\to^* t$ which are proved with respect to an inference system ${\cal I}$ as with Smullyan's \emph{elementary formal systems}~\cite{smullyan1961TheoryFormalSystems}.
%
A \emph{finite proof tree} $T$ in  ${\cal I}$
with root $G$ (an atom) is either:
\begin{itemize}
\item 
%(i) 
an \emph{open goal}, simply denoted as $G$; 
or
\item 
%(ii) 
a \emph{derivation tree} %with root $G$, 
denoted as
$
\frac{T_1 \quad \cdots \quad T_n}{G}(\rho)
$,
where %$G \in \mathit{Form}(\mathcal{S})$,
%$G$ is as above, 
$T_1$,\ldots,$T_n$ are finite proof trees in ${\cal I}$ (for $n\geq 0$; if $n=0$ instead of 
$\frac{}{G}(\rho)$
we just write
$\overline{G}$), and 
$\rho : \frac{B_{1} \cdots  B_{n}}{B}\in{\cal I}$ is an (elementary) inference rule
% from $\GLinferenceOf{\cR}$,
(i.e., $B_1,\ldots,B_n$ and $B$ are atoms) 
such that
%\[
%\frac{\rootTree(T_1) \quad \cdots \quad \rootTree(T_n)}{G}
%\]
$G = \sigma(B)$, and
$root(T_1)=\sigma(B_{1}), \ldots, root(T_n)=\sigma(B_{n})$
for some substitution $\sigma$. % \in\mathit{Sub}(\mathcal{S})$.
\end{itemize}
%
%\begin{remark}[Rules in $\GLinference$ are generic]
%\label{RemGenericityOfInferenceRules}
Note %from the previous statement 
that 
inference rules $\frac{B_{1} \cdots  B_{n}}{B}$
in ${\cal I}$ are viewed
as \emph{schemes of rules} where $B$ should \emph{match} the goal $G$
with matching substitution $\sigma$.
%\end{remark}
%
A finite proof tree $T$ is \emph{closed} if it %whenever it is finite and 
contains no open goals.
%
\begin{definition}[Provable atom]\label{DefProbableGoal}
An atom $G$ is \emph{provable} in ${\cal I}$, %written 
%$\proofInISof{\GLinference}{\GLatom}$,
if there is a closed proof tree $T$ with $root(T)=G$ using ${\cal I}$.
\end{definition}

%
\begin{figure}[t]
\begin{gather*}
(\emph{Rf})\ \frac{}{x \rightarrow^{\ast} x}
\qquad\qquad
(\emph{Co})\ \frac{x \rightarrow y\qquad y \rightarrow^{*} z}{x \rightarrow^{*} z}
\\[0.3cm]
(\emph{Pr})_{f,i}\ \frac{x_i \rightarrow y_i}{f(x_{1},\ldots,x_{i},\ldots,x_{k})\rightarrow f(x_{1},\ldots,y_{i},\ldots,x_{k})}
\\[0.3cm]
(\emph{Rl})_{\ell \to r \Leftarrow s_1\approx t_1,\ldots,s_n\approx t_n}\ \frac{s_1\to^*t_1\quad \cdots \quad s_n\to^*t_n}{\ell \rightarrow r}
\end{gather*}
\caption{Generic Inference rules for conditional rewriting (\emph{oriented} CTRSs)}
\label{FigInferenceRulesForConditionalRewriting}
\end{figure}
%
For \emph{oriented} CTRSs ${\cal R}=({\cal F},R)$, where conditions $u\approx v$ in rules from $R$ are always interpreted as reachability goals $u\to^*v$, such an inference system is built from the generic rules in Figure \ref{FigInferenceRulesForConditionalRewriting} as follows:
\begin{eqnarray*}
{\cal I}({\cal R}) & = & \{(\emph{Rf}),(\emph{Co})\}\cup\{(\emph{Pr})_{f,i}\mid f\in{\cal F}, 1\leq i\leq ar(f)\}\cup\{(\emph{Rl})_\alpha\mid\alpha\in R\}
\end{eqnarray*}
%
Now, we say that $s$ rewrites to $t$ in ${\cal R}$, written $s\to_{\cal R}t$ if there is a well-formed proof tree for the atomic goal $s\to t$ in 
${\cal I}({\cal R})$.

It is well-known that,  in general, it is undecidable whether a \emph{single} conditional rewriting step $s\to_{\cal R} t$ can be issued.
%
In the proof-theoretical description of conditional rewriting above, this has to do with the possibility of having \emph{infinite proof trees} for an atomic goal $s\to t$.
%
From a termination point of view this led to the notion of \emph{operational termination}~\cite{lucas2005OperationalTerminationConditional}. 
%
\begin{definition}
A CTRS ${\cal R}$ is \emph{operationally terminating} if no goal $s\to t$ has an associated infinite proof tree in ${\cal I}({\cal R})$.
\end{definition}
%
The notion of \emph{termination} of a CTRS remains as usual: 
\begin{definition}
A CTRS ${\cal R}$ is \emph{terminating} if there is no infinite sequence $t\to_{\cal R}t'\to_{\cal R}\cdots$ of conditional rewriting steps.
\end{definition}
%
Operational termination implies termination but not vice versa.

\begin{example}
The following (oriented) CTRS
\begin{eqnarray}
a & \to & b \Leftarrow a \approx b
\end{eqnarray}
is \emph{terminating} (the one-step rewrite relation is empty), but it is \emph{not operationally terminating}, as there is an infinite proof tree for $a\to b$ using ${\cal I}({\cal R})$ as above:
\[
\dfrac{\dfrac{\dfrac{\vdots}{a \to b} \quad \overline{b \to^* b}}{a \to^{*} b}}{a \to b}
\]
\end{example}

\section{Competition Categories}

The CTRS formalism is used by the following categories.

\subsection{CTRS Termination}

The CTRS Termination category is concerned with the question 
``Are there infinite sequences of conditional rewriting steps?''.
%
%In other words, does every computation defined by the conditional rewrite rules eventually reach an irreducible term? 

\begin{motivation}
Although guaranteeing operational termination is the most appropriate situation for a practical use of CTRSs, termination is also important to prove other important properties of conditional rewriting, remarkably \emph{confluence}. 
%
The well-known characterization, by Knuth-Bendix and Huet, of confluence of terminating TRSs ${\cal R}$ as the joinability of all its critical pairs  also holds for CTRSs, provided that conditional critical pairs are considered together with the so-called \emph{Conditional Variable Pairs}~\cite{lucas2024LocalConfluenceConditional}.
%
Since termination is \emph{weaker} than operational termination, for the purpose of, e.g., proving confluence as explained above, it is better to try to prove termination than operational termination of CTRSs.
\end{motivation}

\begin{formalproblem}
  {CTRS Termination}
  {CTRS $\Rules$}
  {Does $\Rules$ terminate?}
  {See \hyperref[scoring_termination]{Termination Scoring}}
\end{formalproblem}

\begin{example}
The CTRS $\mathcal{R} = \{a \rightarrow b \Leftarrow f(a) \approx b\}$ has the following input syntax:

\begin{center}
\begin{tabular}{l}
\texttt{(format CTRS oriented)} \\
\texttt{(fun f 1)} \\
\texttt{(fun a 0)} \\
\texttt{(fun b 0)} \\
\texttt{(rule a b (= (f a) b))}
\end{tabular}
\end{center}

and is terminating.
\end{example}

\subsection{CTRS Operational Termination}

The CTRS Operational Termination category is concerned with the question 
``Are there infinite proof trees for conditional rewriting steps?''. 

\begin{motivation}
For the user of a CTRS, operational termination provides important advantages which do \emph{not} follow from termination.
%
In particular, operational termination guarantees decidability of rewriting goals.
\end{motivation}
\begin{formalproblem}
  {CTRS Operational Termination}
  {CTRS $\Rules$}
  {Is $\Rules$ operationally terminating?}
  {See \hyperref[scoring_termination]{Termination Scoring}}
\end{formalproblem}

\begin{example}
The terminating CTRS $\mathcal{R}$ from the previous section, which has the same input syntax, is not operationally terminating, as it can generate an infinite proof tree.
\end{example}
